\newif\ifglobal

%%%%%%%%%%%%%%%%%%%%%%%
%%% Compile to view %%%
%%%%%%%%%%%%%%%%%%%%%%%

%%%%%%%%%%%%%%%%%%%%%%%%%%%%%%%%%%%%%%%%%%%%%%%%%%%
%%% Readme:                                     %%%
%%% https://www.overleaf.com/read/bwdjvfjksvjy/ %%%
%%%%%%%%%%%%%%%%%%%%%%%%%%%%%%%%%%%%%%%%%%%%%%%%%%%

% >>> M?: marks a module setting number or important notice

% >>> M4: Set {./relative-path} to external files for this module <<<
\def \modulefiles {./66-VAD}
% To add an external file, use:
% (Replace '---' with actual filename)
% '\input{\modulefiles/tables/---.tex}' for tables
% '\includegraphics{\modulefiles/figures/---}' for figures
% '\subfile{\modulefiles/---}' for register files


% >>> M1: This module is in global project? <<<
% 'YES' - uncomment; 'NO' - comment out
\globaltrue

\ifglobal
    \documentclass[main\_\_CN.tex]{subfiles}
   
\else
    % Font "FandolSong-Regular" used by ctex does not have GSUB table.
    % It causes warnings that do not affect compilation and can be ignored.
    \PassOptionsToPackage{quiet}{fontspec}

    \documentclass[openany, 10pt]{book}
   
    % >>> M2: In ./00-shared/preamble-trm-module, also find
    %         settings for visibility of todo notes, watermark <<< 
    \usepackage{preamble-trm-module}


    % Enable Chinese typesetting
    \usepackage{ctex}

    % >>> M3: Temporary labels in Chinese
    %         you can add your own labels there <<<
    \myexternaldocument{./00-shared/temp-labels-cn}

\fi %\ifglobal


\setcounter{secnumdepth}{4}

% >>> M5: If this module needs any updates in
% './00-shared/preamble-trm-module.sty'
% (include additional package, etc.), contact your module PM
% or create the file './00-shared/changelog.tex' make the
% necessary updates yourself and document them in this file <<<

% Make text left-aligned, comment out for full justification
\raggedright


\begin{document}


% Sets language into which pre-defined bits of text are translated
\selectlanguage{chinese}

% Do not delete this block
\ifglobal
% Add the following front matter if
% this module is outside of global project
\else
    \InputIfFileExists{readme}{}{}
    {\let\clearpage\relax \listoftodos}
    {\let\clearpage\relax \listofdonetodos}
    \tableofcontents
    %\listoftables
    %\listoffigures
\fi


%%%%%%%%%%%%%%%%%%%%%%%%%
%%% TRM Module Begins %%%
%%%%%%%%%%%%%%%%%%%%%%%%%

\hypertarget{vad}{}
\chapter{语音活动检测 (VAD)}
\label{mod:vad}

\section{概述}
\chipname{} 集成了语音活动检测 (VAD) 模块，硬件实现了语音唤醒等多媒体功能的第一阶段算法。此外，它还为低功耗语音唤醒解决方案提供硬件支持。

\section{主要特性}
\label{sec:vad-features}

VAD 模块支持以下功能：

\begin{itemize}
    \item VAD 算法按帧处理音频数据，每帧 256 个数据点，数据采样率为 8 kHz，位宽 16 位
    \item 2 KB 缓存，可保留多达 4 帧数据
    \item 独立的系统唤醒源
    \item 可配置中断源
    \item 算法参数可灵活配置
\end{itemize}



%%%%%%%%%%%%%%%%%%%%%%%%%%%%%%
%%% Architectural Overview %%%
%%%%%%%%%%%%%%%%%%%%%%%%%%%%%%

\section{系统架构}
\label{sec:vad-arch-overview}

图 \ref{fig:vad-sys-diagram} 为 \chipname{} VAD 模块的结构框图。\chipname{} VAD 模块包含：

\begin{itemize}
    \item 能量阈值检测模块
    \item 快速傅里叶变换 (FFT) 计算模块
    \item 长时频谱散度 (LTSD) 计算模块
\end{itemize}

\begin{figure}[H]
    \centering
    \fbox{\includegraphics[width=0.7\textwidth]{\modulefiles/figures/vad_arch.pdf}}
    \caption{\chipname{} VAD 系统架构}
    \label{fig:vad-sys-diagram}
\end{figure}

VAD 模块从 LP I2S memory 读取 8kHz 采样率，16 位宽的音频数据，通过能量检测、FFT、LTSD 三级运算，检测语音活动状态，并将 VAD 状态信号发送给 LP I2S 模块，用以生成对应中断和唤醒信号。


%%%%%%%%%%%%%%%%%%%%%%%%%%%%%%
%%% Functional Description %%%
%%%%%%%%%%%%%%%%%%%%%%%%%%%%%%

\section{功能描述}
\label{sec:vad-func-descr}
\chipname{} VAD 模块的功能灵活可配置。用户可以通过配置算法参数和唤醒源，满足实际场景需求。

\subsection{算法参数配置}
\label{sec:vad-param-config}


\textbf{初始帧个数}

在 VAD 模块开始运行后，将首先进入初始化阶段，并自动学习噪声参数，初始化阶段的耗时取决于用户配置的初始帧个数。合适的初始帧设置将很好地抑制噪声，显著降低误触发率，而过长的初始帧设置会使得初始帧内的语音活动被遗漏。请根据实际情况，配置合适的初始帧个数。

用户应尽量在初始化阶段使设备处于目标环境噪声中，设定100－200帧长度的初始帧。

用户可以通过配置 \hyperref[fielddesc:LPI2SPARAMINITFRAMENUM]{LP\_I2S\_PARAM\_INIT\_FRAME\_NUM} 来配置初始帧个数。
\begin{tiplisting}
    建议在初始化阶段让VAD模块学习噪声参数，但VAD模块同样会在运行期间不断学习并更新噪声参数。因此，在不关心初始帧中误触发率的场景下，也可以减少初始帧的个数，直接让VAD模块进入运行状态。
\end{tiplisting}


\textbf{能量阈值}

VAD 模块在语音活动检测前会先判定能量阈值。设定能量阈值可以过滤能量过小的语音活动（例如远处的语音活动），避免触发唤醒。设定合适的能量阈值能够降低误触发率，并节省功耗。

VAD 模块定义语音帧的能量值为一帧数据的平方均值，并实时计算一帧数据的平方和，当该平方和小于能量阈值 ×256 时，将不会继续后续计算。用户可以通过 \hyperref[fielddesc:LPI2SPARAMMINENERGY]{LP\_I2S\_PARAM\_MIN\_ENERGY} 来配置能量阈值。



\textbf{通带能量检查}

VAD 通带能量检查判定通带能量在总频域中的占比是否达到要求。在不同的环境中，通带能量检查可能降低误触发率，但也可能增大漏检率。请根据实际环境判断是否需要开启通带能量检查。

用户可以通过 \hyperref[fielddesc:LPI2SPARAMSKIPBANDENERGY]{LP\_I2S\_PARAM\_SKIP\_BAND\_ENERGY} 来开启或关闭通带能量检查。

\begin{itemize}
    \item 0：开启通带能量检查
    \item 1：关闭通带能量检查
\end{itemize}


\textbf{语音活动检测}

VAD 模块的语音活动状态分为语音活动状态和非语音活动状态，内部包含语音帧计数器 (speech\_count) 和非语音帧计数器 (silent\_count) 来控制状态之间的跳转。两种状态的跳转关系如下：

\begin{itemize}
    \item 非语音活动状态下：
    \begin{itemize}
        \item 当前帧有语音活动时，语音帧计数器 (speech\_count) 加一；
        \item 当前帧无语音活动时，语音帧计数器 (speech\_count) 减一（到0为止）；
        \item 当语音帧计数器 (speech\_count) 的值大于 \hyperref[fielddesc:LPI2SPARAMHANGOVERSPEECH]{LP\_I2S\_PARAM\_HANGOVER\_SPEECH} 时，VAD 模块进入语音活动状态，同时归零非语音帧计数器 (silent\_count）；
        \item VAD 模块在非语音活动状态下，会不断学习并更新噪声信息。
    \end{itemize}
    \item 语音活动状态下：
    \begin{itemize}
        \item 当前帧有语音活动时，语音帧计数器 (speech\_count) 加一，同时非语音帧计数器 (silent\_count) 减一（到0为止）；
        \item 当前帧无语音活动时，非语音帧计数器 (silent\_count) 加一，此时：
        \begin{itemize}
            \item 如果语音帧计数器 (speech\_count) 大于 \hyperref[fielddesc:LPI2SPARAMMINSPEECHCOUNT]{LP\_I2S\_PARAM\_MIN\_SPEECH\_COUNT}，且
            \begin{itemize}
                \item 非语音帧计数器 (silent\_count) 大于 \hyperref[fielddesc:LPI2SPARAMHANGOVERSILENT]{LP\_I2S\_PARAM\_HANGOVER\_SILENT}, 则 VAD 模块跳转进入非语音活动状态，同时清零语音帧计数器 (speech\_count) 和非语音帧计数器 (silent\_count)；
                \item 非语音帧计数器 (silent\_count) 小于 \hyperref[fielddesc:LPI2SPARAMHANGOVERSILENT]{LP\_I2S\_PARAM\_HANGOVER\_SILENT}, 则 VAD 模块维持语音活动状态；
            \end{itemize}
            \item 如果语音帧计数器 (speech\_count) 小于等于 \hyperref[fielddesc:LPI2SPARAMMINSPEECHCOUNT]{LP\_I2S\_PARAM\_MIN\_SPEECH\_COUNT}，则 VAD 模块直接跳转进入非语音活动状态，同时清零语音帧计数器 (speech\_count）；
        \end{itemize}
        \item VAD 模块在语音活动状态下，会暂停更新噪声信息，而当语音帧计数器 (speech\_count) 大于 \hyperref[fielddesc:LPI2SPARAMMAXSPEECHCOUNT]{LP\_I2S\_PARAM\_MAX\_SPEECH\_COUNT} 时，视为进入了稳定的新环境中，会继续学习并更新噪声信息。
    \end{itemize}
\end{itemize}


\textbf{状态寄存器}

VAD 模块将一些重要运算变量作为状态寄存器（只读），以方便用户进行调试，状态寄存器包含：
\begin{itemize}
    \item \hyperref[fielddesc:LPI2SVADFLAG]{LP\_I2S\_VAD\_FLAG}：表示当前 VAD 模块的语音活动状态。1代表当前VAD模块处于语音活动状态下，0代表当前VAD模块处于非语音活动状态下。
    \item \hyperref[fielddesc:LPI2SENERGYENOUGH]{LP\_I2S\_ENERGY\_ENOUGH}：表示当前帧的能量阈值检测是否通过。1 代表当前帧通过了能量阈值检测，0 代表当前帧未通过能量阈值检测。
    \item \hyperref[fielddesc:LPI2SSPEECHCOUNTOB]{LP\_I2S\_SPEECH\_COUNT\_OB}：表示语音帧计数器 (speech\_count) 的值。
    \item \hyperref[fielddesc:LPI2SSILENTCOUNTOB]{LP\_I2S\_SILENT\_COUNT\_OB}：表示非语音帧计数器 (silent\_count) 的值。
\end{itemize}


\subsection{唤醒源配置}

当 VAD 模块处于语音活动状态时，会触发唤醒源 LP\_I2S\_WAKEUP。用户可以通过唤醒源配置（请参考章节 \ref{mod:ulp} \textit{\nameref{mod:ulp}}），使 VAD 成为 \chipname{} 的唤醒源之一。


%%%%%%%%%%%%%%%%%%%%%
%%% Interrupts %%%
%%%%%%%%%%%%%%%%%%%%%

\section{中断}

\chipname{} VAD 与 LP I2S 模块共用一个中断信号 LP\_I2S\_INTR。VAD 会生成以下两个中断源：

\begin{itemize}
\item 
\label{int:LPI2SVADDONEINT}LP\_I2S\_VAD\_DONE\_INT：当 VAD 结束一帧数据的运算时，触发此中断。
\item 
\label{int:LPI2SVADRESETDONEINT}LP\_I2S\_VAD\_RESET\_DONE\_INT：当 VAD 完成复位操作时，触发此中断。
\end{itemize}

每个中断源都有一组通用的配置寄存器，详见章节 \hyperref[interrupt-config-registers]{\textit{中断配置寄存器}}。具体寄存器名称请参考章节 \ref{sec:vad-reg-summ} \textit{\nameref{sec:vad-reg-summ}}。


%%%%%%%%%%%%%%%%%%%%%%%%%%%%%%
%%% Programming Procedure  %%%
%%%%%%%%%%%%%%%%%%%%%%%%%%%%%%

\section{配置流程}
\label{sec:vad-prog-procedures}


\subsection{自动运算模式}

\begin{enumerate}
    \item 根据 \ref{sec:vad-param-config} 小节中的描述配置 VAD 参数。
    \item 开启 VAD 自动检测模式，向 \hyperref[fielddesc:LPI2SVADEN]{LP\_I2S\_VAD\_EN} 写 1。
    \item 配置 LP I2S 处于 16 位接收模式，并开始接收数据，接收的数据将存入 LP I2S memory。
    \item VAD 模块会在每次 LP I2S 接收满一帧数据后，自动对该帧数据进行 VAD 运算。
    \item （可选）当 VAD 模块触发唤醒源/中断信号后，LP I2S memory 中已经存储触发 VAD 唤醒的语音帧，和当前语音帧的前3帧语音帧。用户可以通过 LP core 读取 LP I2S memory 中的该4帧语音数据，以进行后续语音信息处理。关于 LP I2S memory 的大小和数据读取方式，请参考章节 \ref{mod:lpi2s} \textit{\nameref{mod:lpi2s}} > 小节 \ref{sec:lpi2s-memory} \textit{\nameref{sec:lpi2s-memory}}。
\end{enumerate}

\subsection{手动运算模式}
\begin{enumerate}
    \item 根据 \ref{sec:vad-param-config} 小节中的描述配置 VAD 参数。
    \item 用户手动向 LP I2S memory 中存储语音数据，存储步骤请参考章节 \ref{mod:lpi2s} \textit{\nameref{mod:lpi2s}} > 小节 \ref{sec:lpi2s-memory} \textit{\nameref{sec:lpi2s-memory}}。
    \item 存储完语音数据后，向 \hyperref[fielddesc:LPI2SVADFORCESTART]{LP\_I2S\_VAD\_FORCE\_START} 写 1， VAD 将会运行一帧数据的运算。
    \item VAD 结束一帧数据的运算时，将生成 \hyperref[int:LPI2SVADDONEINT]{LP\_I2S\_VAD\_DONE\_INT} 中断。
\end{enumerate}


% >>> If your module has no registers, delete sample text
%       for Sections Register Summary and Registers below <<<

%%%%%%%%%%%%%%%%%%%%%%%%
%%% Register Summary %%%
%%%%%%%%%%%%%%%%%%%%%%%%

% >>> Update the sample text below and insert
%     the info on register summary into the subfile <<<

\newpage
\hypertarget{vad-reg-summ}{}
\section{寄存器列表}
\label{sec:vad-reg-summ}

本小节的所有地址均为相对于 VAD 基地址的地址偏移量（相对地址），具体基地址请见章节 \ref{mod:sysmem} \textit{\nameref{mod:sysmem}} 中的表 \ref{tab:sysmem-base-address}。

请查看章节 \hyperref[glossary-access-types]{\textit{寄存器的访问类型}}，了解“访问”列缩写的含义。

\subfile{\modulefiles/66-VAD-reg-summ__CN}


%%%%%%%%%%%%%%%%%
%%% Registers %%%
%%%%%%%%%%%%%%%%%

% >>> Update the sample text below and insert
%      the info on registers into the subfile <<<

\newpage
\section{寄存器}

本小节的所有地址均为相对于 VAD 基地址的地址偏移量（相对地址），具体基地址请见章节 \ref{mod:sysmem} \textit{\nameref{mod:sysmem}} 中的表 \ref{tab:sysmem-base-address}。

关于保留 (reserved) 域的处理，请查看章节 \hyperref[programming-reserved-field]{\textit{如何配置寄存器的保留域}}。

\subfile{\modulefiles/66-VAD-reg__CN}


\end{document}
